% Modernised reading — fonds Alexandre Grothendieck, shelfmark 9, pages 1-28.
%
% This is an INTERPRETATION, not a transcription. It re-reads the pages in
% today's notation and vocabulary, states the hypotheses the manuscript leaves
% implicit, and drops the critical apparatus so that the mathematics reads
% straight through. Everything the brackets and underlines carried moves into a
% footnote. In any dispute of substance the transcription governs: whoever
% cites Grothendieck must cite batch-01.fr.tex and batch-02.fr.tex.
%
% Pass: Opus 5.5 (claude-opus-5-5), 2026-10-03 — first pass, unchecked against
% the pages by a human. Not measured against the Opus 5 reference reading of
% folder 115. The folder was read whole — both batches, pages 1-28 — before
% any of this was written, and it is written as ONE file for the shelfmark.
%
% Notation fixed for the whole folder (spine section): A^* the dual abelian
% scheme and u^* the dual morphism (his star, kept); (-)^∨ the linear dual;
% t_A = Lie(A) (his underlined t), ω_A = det t_A^∨ (as on p. 20; NOT the module
% t_A^∨, which p. 2 also calls ω_G — the module is written t_G^∨ here);
% H^1_dR(A/S) (his DR^1); H(G) kept for the p-divisible case; E(G) for the
% universal vector extension (his G~); A[p], A[p]_ét (his left subscript pA);
% σ for the Frobenius of W(k) (his π); Δ, δ, φ̃, α_A as on pp. 23-28.
%
% Substantive departures from the pages, each footnoted where it occurs:
%   - p. 4-7: A/(F^{n+1},V^{n+1}) is killed by p^{n+1}, so the trace form
%     takes values in W_{n+1}, not W_n as written;
%   - p. 7: « u·V = π(u∘F) » and « Vu = u∘F_g » are slips for V;
%   - p. 8: θ = uθ' is written for θ = θ'∘u;
%   - p. 11: K(W(k'),I) ≃ R_{W(k')}(I) is weakened to a homomorphism;
%   - p. 12: the Z_p(1) counterexample needs inertia (not just G_K) to act
%     non-trivially on μ_p;
%   - p. 13: « égale à n = dim A_K » is read as 2g, the rank (the page's own
%     elliptic case says 2);
%   - p. 19: « t_{A*} → t_A » is a variance slip; the middle component is a
%     section of Hom(t_A, t_{A*}), the tangent map of φ̃ (as the page then says);
%   - p. 23: target ⊗Q_ℓ read ⊗Q_p; det T_p(A^*) read with its « ét »;
%   - p. 24: « u' = φ̃uφ̃^{-1} » is u^* (an endomorphism of A^*); the involution
%     of End(A)⊗Q is the Rosati u ↦ φ̃^{-1}u^*φ̃;
%   - p. 26: the rationality of Δ is argued on the symmetric elements only,
%     which is what the positivity argument needs.
%
% Announced and never established: the identity Δ(u)Δ(u^*) = deg u of p. 24
% (its source is illegible; p. 25, cancelled, cites a « théorème p-adique ou
% th. caractéristique de Weil »); the claim of p. 21 that ω_A ≃ ω_{A^*}
% exchanges the Hasse invariants; c) and d) of p. 28; the « Remarque » of
% p. 17; the end of the cancelled lemma of p. 8. The pencil « faux par Tate »
% of p. 13 is recorded, not reconstructed.
\documentclass[11pt,a4paper]{article}
\input{../preamble/grothendieck}

\folder{9}
\pages{1}{28}
\dating{[vers 1969]}
\watermark{Édition de démonstration\\interprétation personnelle de l'œuvre}
\foldertitle{Groupe de Barsotti-Tate et cristaux (jusqu’en été 1969) : notes manuscrites (s.d.) — lecture modernisée du dossier entier}

\begin{document}

\section*{Résumé}

\begin{resume}
Ces vingt-huit pages parlent de ce qui reste d'une variété abélienne quand on
la regarde « à travers le nombre premier $p$ » en caractéristique $p$. Elles ne
forment pas un mémoire suivi : ce sont cinq séries de feuillets, chacune sous
sa chemise, qui tournent autour d'un même objet et posent sur lui des questions
voisines.

L'objet est le suivant. Une variété abélienne $A$ de dimension $g$ --- une
courbe elliptique si $g = 1$ --- a, pour tout nombre premier $\ell$ différent de
la caractéristique, un « module de Tate » $T_\ell(A)$ : un réseau de rang $2g$
qui enregistre les points de torsion d'ordre une puissance de $\ell$. Pour
$\ell = p$ ce réseau se mutile : il perd la moitié de son rang, ou plus. Ce qui
le remplace est un objet linéaire, le \emph{module de Dieudonné}, muni de deux
opérateurs $F$ et $V$ dont le produit est $p$. Réduit modulo $p$, c'est
un espace vectoriel de dimension $2g$, la cohomologie de De Rham, avec un
sous-espace privilégié de dimension $g$, la filtration de Hodge.

Le problème qui revient d'une série à l'autre est celui-ci : qu'est-ce que
cet objet permet de \emph{canoniser} ? Pour $\ell \neq p$, beaucoup de choses
sont gratuites --- l'accouplement de Weil identifie le module de Tate de $A$
au dual de celui de la variété duale $A^*$. Pour $\ell = p$, rien ne l'est.

L'idée qui arrive est que tout devient possible dans le cas \emph{ordinaire}.
Une image : un espace muni d'un opérateur peut avoir une partie où
l'opérateur est inversible et une partie où il est nilpotent ; quand ces
deux parties remplissent l'espace et se séparent proprement, tout ce qui est
défini sur l'une se transporte. Une variété abélienne ordinaire est celle où
$F$ et $V$ découpent ainsi le module de Dieudonné en deux moitiés de même
taille, l'une « étale », l'autre « multiplicative ». Alors, montre le
dossier, la cohomologie de De Rham se scinde canoniquement, chaque moitié
devient un système local étale tensorisé par les fonctions, et un
« transport parallèle » canonique existe sur l'espace des déformations ---
ce qui, en caractéristique $p$, n'arrive \emph{que} dans le cas ordinaire.
L'obstruction a un nom : la section qui s'annule exactement là où la variété
cesse d'être ordinaire, et que les pages baptisent « invariant de Hasse ».
Les dernières pages en tirent un isomorphisme canonique entre les
déterminants des parties étales des modules de Tate $p$-adiques de $A$ et de
$A^*$ : on le fabrique à partir d'une polarisation, on divise par la racine
carrée de son degré, et l'on montre --- par un argument de positivité d'une
grande économie --- que le résultat ne dépend pas de la polarisation choisie.

Entre deux, une série dactylographiée pose une question d'une autre nature :
peut-on attacher à un groupe $p$-divisible sur un corps local un
« caractère » de l'inertie, comme on le fait pour $\ell \neq p$, et ce
caractère est-il lu par la seule $p$-torsion ? La réponse est non, et la
page le démontre elle-même ; la question suivante porte, au crayon, la mention
« faux par Tate ». Ce dossier montre ainsi le travail en train de se faire :
un lemme écrit puis barré d'un « faux », une page entière annulée, des
questions laissées pour qui voudra.

Les noms modernes à chercher ensuite : la théorie de Dieudonné et l'extension
vectorielle universelle, la théorie de Serre--Tate des variétés abéliennes
ordinaires, la filtration conjuguée, l'invariant de Hasse, et --- pour la
question des caractères --- la réduction potentiellement bonne et les
caractères de Brauer.

\keywords{Dieudonné module, p-divisible group, universal vector extension, Cartier duality, Gauss-Manin connection, potentially good reduction, Brauer character, Hodge filtration, conjugate filtration, ordinary abelian variety, Serre-Tate theory, Hasse invariant, Rosati involution}
\end{resume}

\pagerange{1}{28}
\section*{Le fil du dossier, et ses conventions}

Le dossier n'a pas une colonne vertébrale unique ; il a un objet commun. Cinq
séries de feuillets, non paginées par lui, séparées par des chemises dont
les titres sont repris ici :

\begin{enumerate}
\item \emph{Pages 1 à 7} (chemise : « Dualité des modules de Dieudonné de
groupes de B.T.\ duaux l'un de l'autre ») --- la réduction modulo $p$ du
module de Dieudonné, sa dualité, et l'anneau de Dieudonné tronqué muni d'une
forme trace.
\item \emph{Page 8}, annulée --- une question sur la connexion de
Gauss--Manin et un lemme sur les torseurs d'homotopies.
\item \emph{Pages 9 à 13} (chemise : « Caractères des groupes de B.T.\ à
réduction semi-stable virtuelle essentielle ») --- un tapuscrit sur le
caractère d'un groupe $p$-divisible sur un corps local.
\item \emph{Pages 14 à 20} (chemise : « Détermination des
“stratifications canoniques” sur les déformations ») --- les déformations
d'une variété abélienne, la nécessité de l'ordinarité, la description étale,
la polarisation, l'invariant de Hasse.
\item \emph{Pages 21 à 28}, à la suite sans chemise mais annoncées par celle
de la page 14 (« L'isom.\ de nature globale : $\det T_p(A)_{\text{ét}} \simeq
\det T_p(A^*)_{\text{ét}}$ ») --- la construction de cet isomorphisme.
\end{enumerate}

Les séries 1, 4 et 5 se tiennent : elles parlent du même espace $H$ muni de
$F$ et $V$, d'abord pour un groupe $p$-divisible sur un corps (page 2), puis
pour une variété abélienne sur une base. Les séries 2 et 3 sont des
digressions, la seconde d'une autre main (une machine à écrire) et d'un
autre côté du sujet (les représentations galoisiennes).

\emph{Conventions, valables pour tout le dossier.} $k$ est un corps parfait de
caractéristique $p > 0$, $W = W(k)$ ses vecteurs de Witt, $W_n = W/p^n W$, et
$\sigma$ le Frobenius de $W$.\footnote{Il le note $\pi$, aux pages 4, 7 et
20 ; à la page 20 la même lettre désigne une flèche $t_A^{(p)} \to t_A$, que
l'on note ici $V$.} Pour un groupe $p$-divisible $G$ (un « groupe de B.T. »),
$G^*$ est le dual de Cartier ; pour un schéma abélien $A$ sur $S$, $A^*$ est le
schéma abélien dual et $u^* : B^* \to A^*$ le dual d'un morphisme $u : A \to
B$ --- l'étoile est la sienne et on la garde. Le signe ${}^\vee$ désigne
toujours un dual \emph{linéaire}. $t_A = \operatorname{Lie}(A)$ est l'espace
tangent à l'origine (son $t$ souligné), $t_A^\vee$ le module des formes
invariantes, et $\omega_A = \det t_A^\vee = \bigwedge^g t_A^\vee$ le
fibré de Hodge, au sens de la page 20.\footnote{À la page 2 il écrit
$\omega_G$ pour le module $t_G^\vee$ lui-même, non pour son déterminant. On
écrit ici $t_G^\vee$ dans ce cas, pour que $\omega$ ait un seul sens.}
$H^1_{\mathrm{dR}}(A/S)$ est ce qu'il note $\mathrm{DR}^1(A/S)$ ; $A[p]$ et
$A[p]_{\text{ét}}$ sont ses ${}_pA$ et $({}_pA)_{\text{ét}}$ ; $M^{(p)}$ est le
tiré en arrière par le Frobenius de la base.

\emph{Ce que le dossier annonce et n'établit pas} : l'identité $\Delta(u)
\Delta(u^*) = \deg u$ sur laquelle repose la page 24 (sa justification est
illisible, et la page 25 qui en esquissait une est annulée) ; la
compatibilité de l'isomorphisme $\omega_A \simeq \omega_{A^*}$ avec les
invariants de Hasse (page 21) ; les points c) et d) de la page 28 ; la
fin du lemme de la page 8. On les retrouve à leur place.

\pagerange{1}{7}
\section*{Le module de Dieudonné réduit et sa dualité (pages 1 à 7)}

\pagerange{2}{2}
\subsection*{Les deux suites exactes (page 2)}

Soit $G$ un groupe $p$-divisible sur $k$, de dimension $g$, dont le dual
$G^*$ est de dimension $g^*$ ; la hauteur est $g + g^*$. L'extension
vectorielle universelle de $G$ est une suite exacte
\[
0 \to \mathbb{V}(\omega_{G^*}) \to E(G) \to G \to 0 ,
\]
où $\mathbb{V}(\omega_{G^*})$ est le groupe vectoriel dont les points sont les
sections de $t_{G^*}^\vee$.\footnote{Il écrit $\tilde G$ pour $E(G)$, et
$\mathbb{V}(\underline{t}_{G^*})$ pour le noyau : dans la convention d'EGA,
$\mathbb{V}(\mathcal{E}) = \operatorname{Spec}\operatorname{Sym}\mathcal{E}$, dont les points sont les formes
linéaires sur $\mathcal{E}$, si bien que les deux écritures désignent le même
groupe. Le nom d'« extension vectorielle universelle » et la notation $E(G)$
sont postérieurs (Mazur et Messing, 1974, à la suite des leçons de
Grothendieck à Montréal en 1970) ; l'objet est déjà là.} En passant aux
algèbres de Lie, on obtient
\[
0 \to t_{G^*}^\vee \to \operatorname{Lie} E(G) \to t_G \to 0 .
\]
Posons $H(G) = \operatorname{Lie} E(G^*)$ --- la page le dit en note, « NB
$H(G) = \underline{t}_{\tilde G^*}$ ».\footnote{L'indice est en partie
recouvert d'une tache ; la lecture $\tilde G^*$ est celle de la transcription,
appuyée sur la même égalité, lisible, quelques lignes plus bas.} C'est
l'analogue, pour $G$, de $H^1_{\mathrm{dR}}$ d'une variété abélienne, et c'est
la réduction modulo $p$ du module de Dieudonné de $G$ --- le « module de
Dieudonné réduit $M \otimes_W k$ » de la chemise. Il s'insère dans
\[
0 \to t_G^\vee \to H(G) \to t_{G^*} \to 0 ,
\qquad
\dim = g,\ g + g^*,\ g^* ,
\]
et porte deux applications, $F : H(G)^{(p)} \to H(G)$ et $V : H(G) \to
H(G)^{(p)}$, de composés nuls. La suite exacte pour $G^*$, c'est-à-dire pour
$H(G^*) = \operatorname{Lie} E(G)$, est la \emph{transposée} de celle de $G$ :
la dualité $H(G^*) \simeq H(G)^\vee$ échange les deux suites, et échange $F$
et $V$.\footnote{Les deux diagrammes sont côte à côte sur la page, avec
« dual de » au-dessus du trait qui les sépare ; la phrase (1) qui énonce
la transposition est à moitié illisible, et l'énoncé précis --- dualité,
échange de $F$ et $V$ --- est notre reconstruction. C'est l'énoncé connu.}

L'énoncé (2) de la page est exact et se dit ainsi :
\[
\operatorname{Ker} F = \operatorname{Im} V = (t_G^\vee)^{(p)} ,
\qquad
\operatorname{Ker} V = \operatorname{Im} F .
\]
Ce second sous-espace, de dimension $g^*$, est l'image de $t_{G^*}^{(p)}$ par
$F$ ; on l'appelle aujourd'hui la \emph{filtration conjuguée}, par opposition à
la filtration de Hodge $t_G^\vee \subset H(G)$. Il n'est pas en général un
supplémentaire de $t_G^\vee$ : c'est exactement ce que mesure l'ordinarité, aux
pages 15 et 21.\footnote{La page laisse en blanc ce que vaut $\operatorname{Ker}
V$ (« ce que ça \ldots », suivi d'une parenthèse illisible). Pour une variété
abélienne, c'est le théorème d'Oda (1969) sur $H^1_{\mathrm{dR}}$ en
caractéristique $p$ ; la page n'y renvoie pas.}

Vient enfin une trichotomie, que l'on énonce en termes de la décomposition de
$G$ en parties connexe et étale :
\begin{itemize}
\item $G$ est connexe (« purement infinitésimal »), c'est-à-dire $G^*$ n'a pas
de partie de type multiplicatif (« sans tore »), si et seulement si $F$ est
nilpotent sur $H(G)$ ;
\item $G^*$ est connexe, c'est-à-dire $G$ n'a pas de partie de type
multiplicatif, si et seulement si $V$ est nilpotent ;
\item $G$ et $G^*$ sont tous deux connexes --- $G$ est « local-local » --- si
et seulement si $F$ et $V$ sont nilpotents.
\end{itemize}
On le vérifie sur les cas extrêmes : pour $G$ étale, $H(G) = t_{G^*}$ est
l'algèbre de Lie d'un tore et $F$ y est bijectif ; pour $G = \mu_{p^\infty}$,
$H(G) = t_G^\vee$ et $F$ y est nul.\footnote{Ces trois lignes sont d'une
lecture presque entièrement incertaine sur la page ; seuls les signes
$\iff$ et la structure en trois cas sont sûrs. Les énoncés donnés sont ceux
que la théorie de Dieudonné rend vrais, et que la lecture de la transcription
suggère. Le point (3), sur la composée $t_{G^*}^{(p)} \to H(G) \to t_{G^*}$ ---
l'application de Hasse--Witt de $G^*$ --- et sur un « Frobenius dans
$H^2(G, \mathbb{G}_m)$ », est trop lacunaire pour qu'on en tire un énoncé.}

\pagerange{4}{7}
\subsection*{L'anneau de Dieudonné tronqué et sa forme trace (pages 4, 6 et 7)}

Soit $\mathcal{A} = W[F, V]$ l'anneau de Dieudonné : $F\lambda =
\lambda^\sigma F$, $V\lambda^\sigma = \lambda V$, $FV = VF = p$. Il est libre
sur $W$, de base les $F^i$ ($i \geq 0$) et les $V^j$ ($j \geq 1$). Pour $n
\geq 0$, on pose
\[
\mathcal{A}_n = \mathcal{A}/(F^{n+1}, V^{n+1}) .
\]
Comme $V^a F^{n+1} = p^a F^{n+1-a}$, l'idéal bilatère contient $p^{n+1}$, et
$\mathcal{A}_n = \bigoplus_{i=0}^n W_{n+1-i} F^i \oplus \bigoplus_{j=1}^n
W_{n+1-j} V^j$.\footnote{La page écrit partout $W_n$ comme anneau de
coefficients de $\mathcal{A}_n$. Avec la troncature choisie, $\mathcal{A}_n$
est tué par $p^{n+1}$ et non par $p^n$ ; on corrige donc l'indice en $W_{n+1}$.
On aurait aussi bien pu garder $W_n$ et tronquer en $F^n$, $V^n$, ce que la
page 6 suggère d'ailleurs (« si $M$ tué par $F^n$, $V^n$ »).}

Il existe une unique forme $W$-linéaire $\varphi_n : \mathcal{A}_n \to W_{n+1}$
telle que $\varphi_n(1) = 1$ et $\varphi_n(F^i) = \varphi_n(V^i) = 0$ pour $i
\geq 1$ ; elle est linéaire à gauche et à droite sur $W$, et vérifie
\[
\varphi_n(Fz) = \varphi_n(zF)^\sigma , \qquad \varphi_n(zV) =
\varphi_n(Vz)^\sigma .
\]
On le vérifie sur la base : pour $z = \lambda V$, $Fz = p\lambda^\sigma$ et $zF
= p\lambda$. La page le fait, pour la linéarité, en écrivant $z = a + \sum
\lambda_i F^i + \sum V^i \mu_i$, et s'arrête au moment de traiter $F$.

La forme $\langle x, y \rangle = \varphi_n(xy)$ est le point de la page 7. Ses
règles sont
\[
\langle \lambda x, y \rangle = \lambda \langle x, y \rangle , \quad
\langle x, y \lambda \rangle = \langle \lambda x, y \rangle , \quad
\langle x, yF \rangle = \sigma^{-1}\langle Fx, y \rangle , \quad
\langle x, yV \rangle = \sigma\langle Vx, y \rangle ,
\]
et $\langle x, ay \rangle = \langle xa, y \rangle$ pour tout $a \in
\mathcal{A}_n$, ce qui la factorise par $\mathcal{A}_n \otimes_{\mathcal{A}_n}
\mathcal{A}_n$. Elle est \emph{parfaite} : $F^i$ s'apparie avec $V^i$ en $p^i$,
qui engendre $p^i W_{n+1} \simeq W_{n+1-i}$, et tous les autres produits de
vecteurs de base ont une trace nulle. Il en résulte un isomorphisme de
$\mathcal{A}_n$-bimodules
\[
\mathcal{A}_n \xrightarrow{\ \sim\ } \operatorname{Hom}_W(\mathcal{A}_n,
W_{n+1}) ,
\]
que la page pose avec un point d'interrogation au-dessus de la flèche et
dont elle vérifie la linéarité à droite et à gauche par les deux groupes de
règles ci-dessus.\footnote{Le caractère parfait de l'accouplement est notre
vérification ; la page n'en dit rien, et c'est ce qui fait de la flèche un
isomorphisme.} La structure de module à droite sur $D'(M) =
\operatorname{Hom}_W(M, W_{n+1})$, pour un $\mathcal{A}$-module à gauche $M$
tué par $F^{n+1}$ et $V^{n+1}$, est la transposée :
\[
u \cdot F = \sigma^{-1} \circ u \circ F , \qquad u \cdot V = \sigma \circ u
\circ V , \qquad u \cdot \lambda = u \circ \lambda ,
\]
et l'on a bien $u \cdot F \cdot V = u \cdot V \cdot F = pu$.\footnote{La page
écrit « $u \cdot V = \pi(u \circ F)$ », et plus bas « $Vu = u \circ F_g$ » :
deux lapsus pour $V$, que le calcul $u \cdot F \cdot V = pu$ écrit juste à
côté exige.}

Cela met en place deux foncteurs dualisants pour les modules de Dieudonné de
longueur finie : $D(M) = \operatorname{Hom}_{\mathcal{A}}(M, I)$ où $I =
\varinjlim I_n$, $I_n = \mathcal{A}_n$ vu comme bimodule et les transitions
induites par la multiplication par $p$ ; et $D'(M) = \operatorname{Hom}_W(M,
W_{n+1})$ muni de la structure transposée. L'isomorphisme ci-dessus est ce qui
les identifie. Tous deux changent un module à gauche en module à droite ; on
revient aux modules à gauche par l'anti-involution $s$ de $\mathcal{A}$, triviale
sur $W$, qui échange $F$ et $V$.\footnote{La page dit « involution » ; c'est
un anti-automorphisme ($s(ab) = s(b)s(a)$), et c'est ce qu'il faut pour
changer de côté. On vérifie qu'il est bien défini : $s(\lambda^\sigma F) =
V\lambda^\sigma = \lambda V = s(F\lambda)$.} C'est, sous sa forme d'anneau, la
dualité de Dieudonné qui correspond à la dualité de Cartier ; aujourd'hui on
l'écrit d'ordinaire directement $M \mapsto \operatorname{Hom}_W(M, W[1/p]/W)$,
$F$ et $V$ transposés et échangés.

\pagerange{8}{8}
\section*{Une question sur la connexion de Gauss--Manin (page 8)}

La page est annulée de trois traits, mais lisible.\footnote{On la donne parce
qu'elle se lit et qu'elle situe le dossier ; son annulation interdit d'y voir
autre chose qu'un essai.} Elle demande comment définir \emph{directement} une
connexion intégrable sur $R f_* \Omega^\bullet_{X/S}$ --- ce qu'on appelle la
connexion de Gauss--Manin, construite par Katz et Oda en 1968 et, du point de
vue cristallin, déduite de ce que la cohomologie de De Rham est un cristal.

Le lemme qui suit fournit l'outil. Soient $X$ un topos, $K^\bullet$ et
$L^\bullet$ des complexes de faisceaux abéliens, $\mathcal{G}$ un faisceau
abélien muni de $i : \mathcal{G} \to \underline{\mathrm{Hom}}^{-1}(K^\bullet,
L^\bullet) = \prod_j \underline{\mathrm{Hom}}(K^j, L^{j-1})$, et $P$ un
$\mathcal{G}$-torseur muni de $\theta : P \to Z^0
\underline{\mathrm{Hom}}^\bullet(K^\bullet, L^\bullet)$ --- à chaque section de
$P$ un morphisme de complexes --- tel que
\[
\theta(g \cdot p) = \theta(p) + d\, i(g) + i(g)\, d .
\]
Autrement dit : localement on dispose d'un morphisme de complexes, et deux
choix diffèrent d'une homotopie fournie par $\mathcal{G}$, de façon
cohérente. Alors on construit un morphisme $\varphi : K^\bullet \to L^\bullet$
dans $D(X)$, fonctoriel en $L^\bullet$ et en la donnée $(\mathcal{G}, P, i,
\theta)$ (pour $u : \mathcal{G} \to \mathcal{G}'$, $P \to P'$ compatible, $i =
i' u$, $\theta = \theta' \circ u$) ;\footnote{La page écrit « $\theta = u
\theta'$ », ce qui ne compose pas : $\theta$ part de $P$ et $u$ va de $P$ vers
$P'$.} et la page s'arrête sur « 3°) Si $P$ admet une section, alors » --- on
attend : $\varphi$ est la classe de $\theta(p)$ pour n'importe quelle section
$p$.\footnote{La conclusion est notre complément. Le lemme n'est pas
démontré, et il n'est pas trivial : les morphismes de $D(X)$ ne se recollent
pas en général, et c'est la cohérence des homotopies --- un groupe qui opère,
non une simple famille --- qui permet de conclure.}

\pagerange{9}{13}
\section*{Le caractère d'un groupe de Barsotti--Tate sur un corps local (pages 9 à 13)}

Les pages 10 à 13 sont un tapuscrit, tapé en bleu, corrigé à l'encre et au
crayon de sa main ; son titre est « Caractère d'un groupe de B.T.\ sur un
corps local ».\footnote{Le tapuscrit n'est pas identifié comme publié. La
chemise (page 9) porte un titre plus long, « à réduction semi-stable
virtuelle essentielle », c'est-à-dire, en termes d'aujourd'hui,
potentiellement semi-stable.}

\pagerange{10}{11}
\subsection*{La condition de monodromie et l'homomorphisme trace}

Soit $V$ un anneau de valuation discrète complet, de corps résiduel $k$
parfait de caractéristique $p$, de corps des fractions $K$, et $S =
\operatorname{Spec} V$.\footnote{Le tapuscrit hésite sur le corps
résiduel : « algébriquement clos » est barré au profit de « parfait », ajouté
à l'encre ; la parenthèse « quitte à passer au complété de l'hensélisé
strict » est barrée puis peut-être rétablie. $S$ est employé sans être
introduit.} Un groupe $p$-divisible $M$ sur $K$ \emph{satisfait la condition
de monodromie} s'il admet une filtration par des sous-groupes $p$-divisibles
dont les quotients successifs ont \emph{potentiellement bonne réduction} :
chacun, après une extension finie séparable $K'$ de $K$, s'étend en un
groupe $p$-divisible sur le normalisé $S'$ de $S$ dans $K'$. On peut prendre
$K'$ commun à tous les quotients, galoisien de groupe $G$ : c'est un
\emph{corps de bonne réduction potentielle} pour $M$.

Le prolongement est unique à isomorphisme unique près par le théorème de Tate
(pleine fidélité du passage à la fibre générique), et $G$ opère donc sur lui,
semi-linéairement par rapport à son action sur $S'$.\footnote{Le tapuscrit
précise : démontré « pour l'instant seulement dans le cas d'inégales
caractéristiques » --- c'est le théorème de Tate de 1967, pour $K$ de
caractéristique nulle. Le cas d'égale caractéristique a été établi par de
Jong en 1998.} De là une chaîne d'homomorphismes de groupes de Grothendieck :
\[
K_{BT}(K; K') \to K_{BT}(S'; G) \to K_{BT}(s'; G) \to K(W(k'), G) \to
K(W(k'), I) ,
\]
où $s'$ est le point fermé de $S'$, $k'$ son corps résiduel, la troisième
flèche est le module de Dieudonné (foncteur covariant, « et même vers les
cristaux de $k'$ », dit le crayon), et $I \subset G$ est le sous-groupe
d'inertie, qui opère $W(k')$-linéairement. La dernière étape est la trace :
\[
K(W(k'), I) \to \operatorname{Cent}(I, W(k'))^{G} ,
\]
à valeurs dans les fonctions centrales sur $I$, invariantes sous l'action de
$G$ (par conjugaison sur $I$ et par son action sur $W(k')$).\footnote{La page
écrit $K(W(k'), I) \simeq R_{W(k')}(I)^{G/I}$. Pour un $I$ d'ordre divisible
par $p$, le groupe de Grothendieck des réseaux n'est pas évidemment l'anneau des
représentations ; on n'utilise que l'homomorphisme vers les représentations
de $I$ sur le corps des fractions de $W(k')$, par lequel la trace se
factorise. C'est tout ce que la définition du caractère demande.} L'image de
la classe de $M$ est le \emph{caractère} de $M$ ; il ne dépend pas de $K'$ et
se voit comme une fonction centrale sur le groupe d'inertie (infini) de
$\operatorname{Gal}(\bar K/K)$, qui se factorise par un quotient fini.

\pagerange{11}{12}
\subsection*{Le caractère n'est pas lu sur la $p$-torsion}

Si $M_{K'}$ est \emph{constant}, son module de Dieudonné est $T_p(M) \otimes
W(k')$, et son caractère coïncide, sur les éléments de $I$ d'ordre premier à
$p$, avec le \emph{caractère de Brauer} (le « caractère modulaire ») de la
représentation de $I$ sur $M[p] = M(0)$.\footnote{Ses $M(\nu)$ sont les
$M[p^{\nu+1}]$ : $M(0) = {}_pM$. Le point d'interrogation qui clôt la phrase
est repassé à l'encre ; l'énoncé est vrai, pour la raison donnée.}

Le tapuscrit se demande alors si cela vaut pour tout $M$ satisfaisant la
condition de monodromie, et répond \emph{non}, par un argument propre. Si le
caractère était toujours déterminé par $M[p]$, la théorie de Brauer des
relèvements virtuels forcerait l'égalité avec le caractère de Brauer de $M[p]$
pour tout $M$ ; en particulier, pour $M$ à bonne réduction --- dont le
caractère est constant, égal à la hauteur --- le caractère de Brauer de $M[p]$
serait trivial. Or $M = \mathbb{Z}_p(1)$, c'est-à-dire $\mu_{p^\infty}$, a bonne
réduction et $M[p] = \mu_p$, sur lequel l'inertie opère par le caractère
cyclotomique modulo $p$.\footnote{La page dit « dans le cas où $K$ ne contient
pas les racines $p$-ièmes de $1$ ». L'hypothèse exacte est que l'\emph{inertie}
opère non trivialement sur $\mu_p$, c'est-à-dire que $\mu_p$ n'est pas dans
l'extension non ramifiée maximale de $K$ ; elle équivaut à celle de la page
quand $k$ est algébriquement clos, et elle est satisfaite pour $K =
\mathbb{Q}_p$. Pour $K = \mathbb{Q}_p(p^{1/(p-1)})$, $p$ impair, $K$ ne contient
pas $\mu_p$ mais $K(\mu_p)/K$ est non ramifiée, et le contre-exemple tombe.}

Si $K$ contient $\mu_p$, au contraire, rien ne s'oppose à l'énoncé pour les
bonnes réductions potentielles constantes ou multiplicatives ; reste le cas
où la réduction est connexe à dual connexe (« ind-unipotente »). Une note au
crayon en tire que le caractère serait alors à valeurs dans $\mathbb{Z}_p
\subset W(k)$.

\pagerange{12}{13}
\subsection*{Le cas test des courbes elliptiques}

D'où le cas test. Soient $E_1$, $E_2$ deux courbes elliptiques sur $K$ à bonne
réduction potentielle, avec $E_1[p] \simeq E_2[p]$ et $\mu_p \subset K$. Les
traces $\phi_{E_1}$ et $\phi_{E_2}$ sur l'inertie --- que l'on peut aussi
calculer sur $T_\ell$ pour n'importe quel $\ell$\footnote{Parce que l'inertie
opère sur la réduction $E_{s'}$ par automorphismes, et que la trace d'un
endomorphisme d'une variété abélienne est la même sur $T_\ell$, $\ell \neq p$,
et sur le module de Dieudonné. C'est le « comme on sait » de la page.} ---
coïncident-elles sur les éléments d'ordre premier à $p$ ? Seul le cas d'une
réduction supersingulière (« invariant de Hasse nul ») est intéressant.
Variante : si $A_K[p]$ est constant, la trace vaut-elle $2g$ sur les éléments
de $I$ d'ordre premier à $p$ ?\footnote{Le tapuscrit porte « égale à $n = \dim
A_K$ », tapé au-dessus de mots surfrappés. La trace de l'identité est le rang
$2g$, et le paragraphe suivant, pour la courbe elliptique, donne bien $2$ ; on
corrige.} Il faudrait, dit la page, un exemple où le critère de réduction
semi-stable de Raynaud --- $A[\ell]$ rationnel, $\ell \geq 3$, entraîne
réduction semi-stable --- est en défaut pour $\ell = p$.

Pour une courbe elliptique, en caractéristique résiduelle $\neq 2, 3$, avec
$K'$ minimal, l'inertie de $K'/K$ s'injecte dans $\operatorname{Aut}(E_{s'})$,
groupe cyclique d'ordre $2$, $4$ ou $6$, donc d'ordre premier à $p$ ; la
question revient à savoir si $\phi_E$ est constante égale à $2$ lorsque $E[p]$
est constant. En marge, au crayon : \emph{faux par Tate}.\footnote{Rien sur la
page ne dit ce que Tate objectait, et on ne le reconstitue pas. On note
seulement qu'un automorphisme non trivial d'ordre $2$, $3$, $4$ ou $6$ a pour
trace $-2$, $-1$, $0$ ou $1$, de sorte que la question équivaut à demander si
la constance de $E[p]$ force l'inertie à opérer trivialement sur $E_{s'}$.}

\pagerange{14}{18}
\section*{Les déformations d'une variété abélienne et leurs stratifications (pages 14 à 18)}

\pagerange{15}{15}
\subsection*{La question, et pourquoi elle force l'ordinarité}

Soit $A_0$ une variété abélienne de dimension $g$ sur $k$. La question
(page 15) : quand peut-on associer à toute déformation infinitésimale
$A/\Lambda$ de $A_0$ --- $\Lambda$ une $k$-algèbre locale artinienne de corps
résiduel $k$ --- une \emph{stratification} relative à $k$ sur
$H^1_{\mathrm{dR}}(A/\Lambda)$, compatible au changement de base et
fonctorielle en $A$ ? Une stratification est un transport parallèle à tous les
ordres : un isomorphisme entre les deux images inverses sur $\operatorname{Spf}
(\Lambda \hat\otimes_k \Lambda)$, satisfaisant la condition de cocycle. En
caractéristique $p$, c'est beaucoup plus qu'une connexion : la connexion de
Gauss--Manin existe toujours, une stratification non.

Soit $\mathcal{M}$ l'espace des modules formels de $A_0$ --- lisse sur $k$, de
dimension $g^2$ --- et $A_{\mathcal{M}}$ la déformation universelle. Il faudrait
que $\mathcal{H} = H^1_{\mathrm{dR}}(A_{\mathcal{M}}/\mathcal{M})$ porte une
stratification compatible avec $F : \mathcal{H}^{(p)} \to \mathcal{H}$ et $V :
\mathcal{H} \to \mathcal{H}^{(p)}$. Alors le type de $(\mathcal{H}_x, F_x, V_x)$
serait le même en tout point $x$ que celui de $(\mathcal{H}_0, F_0, V_0)$ ---
c'est-à-dire, en termes d'aujourd'hui, que le type d'Ekedahl--Oort de $A_x[p]$
serait constant sur $\mathcal{M}$. Comme le lieu ordinaire est dense dans
$\mathcal{M}$, cela force $A_0$ à être \emph{ordinaire}.\footnote{La page écrit
« Cela implique sans doute que $A_0$ est “ordinaire” » et ne justifie pas.
L'argument par la densité du lieu ordinaire dans l'espace des déformations
est le nôtre ; ce fait n'est pas sur la page.}

Ordinaire signifie : dans la suite de Hodge
\[
0 \to H^{10} \to H \to H^{01} \to 0 , \qquad H^{10} = t_{A}^\vee =
H^0(A, \Omega^1), \quad H^{01} = t_{A^*} = H^1(A, \mathcal{O}_A) ,
\]
où $\operatorname{Ker} F = \operatorname{Im} V = H^{10(p)}$ et
$\operatorname{Ker} V = \operatorname{Im} F$ comme à la page 2, l'image $\operatorname{Im} F$ est un
supplémentaire de $H^{10}$ ; de façon équivalente, la composée $H^{01(p)} \to
H \to H^{01}$ --- l'opération de Frobenius sur $H^1(A, \mathcal{O}_A)$,
l'application de Hasse--Witt --- est bijective.

\pagerange{15}{17}
\subsection*{La filtration de Hodge est horizontale}

Supposons $A_0$ ordinaire. La stratification, compatible avec $F$, laisse
stable $\operatorname{Im} F = F(\mathcal{H}^{(p)})$, qui est un supplémentaire
de $H^{10}$ ; de même elle laisse stable $\operatorname{Ker} F = H^{10(p)}$.
Cela n'entraîne pas \emph{a priori} que $H^{10}$ soit stable --- c'est
l'« Attention » de la page 16. Ce qu'on obtient : sur $T = S \times S$, $S =
\operatorname{Spf}\Lambda$, soient $E_1$, $E_2$ les deux images inverses de
$H$, $F_1$, $F_2$ celles de $H^{10}$, et $u : E_1 \xrightarrow{\sim} E_2$
l'isomorphisme de stratification ; alors $u(F_1)^{(p)} = F_2^{(p)}$. Les deux
sections $u(F_1)$ et $F_2$ de la grassmannienne de $E_2$ deviennent égales
après changement de base par le Frobenius absolu de $T$. Si $T$ est
formellement lisse sur $k$ --- c'est le cas pour $S = \mathcal{M}$ ---, le
Frobenius est fidèlement plat (théorème de Kunz), donc $u(F_1) = F_2$ ; et le
cas général s'en déduit par fonctorialité, en remontant à l'espace de
modules.\footnote{La page dit « c'est loc.\ fid.\ plat » sans nommer de
théorème. Le théorème de Kunz (1969) caractérise les anneaux réguliers par la
platitude du Frobenius ; pour un anneau de séries formelles sur un corps
parfait, la platitude se voit directement.}

Donc $H^{10}$ est horizontal, et $H$ admet une décomposition fonctorielle,
compatible au changement de base et à la stratification :
\[
H^1_{\mathrm{dR}}(A/S) \simeq t_A^\vee \oplus t_{A^*} ,
\]
le second facteur étant identifié à $\operatorname{Im} F$. Le problème se
ramène à trouver une stratification fonctorielle sur $t_A$, le cas de
$t_{A^*}$ étant celui de $A^*$. La page hésite entre une stratification
« a priori » venue des puissances divisées et une définition directe par les
puissances $p$-ièmes ; « à préciser », dit la marge. La page 18 tranche,
comme on va voir.

La \emph{Remarque} qui clôt la page 17 entrevoit que la position de la
filtration de Hodge par rapport à une structure horizontale sur la base des
déformations devrait donner un principe de classification des déformations en
caractéristique $p$.\footnote{La remarque est presque entièrement d'une
lecture incertaine, et la parenthèse n'est pas fermée. Lue littéralement pour
une stratification, elle serait vide dans le cas ordinaire, puisque la
filtration y est horizontale ; on la lit donc, sous notre responsabilité, comme
visant la structure cristalline. C'est le théorème de Grothendieck--Messing :
les déformations d'un groupe de Barsotti--Tate le long d'un épaississement à
puissances divisées sont classées par les relèvements de la filtration de
Hodge dans le cristal de Dieudonné --- exposé par Grothendieck au congrès de
Nice en 1970 sous un titre qui est celui même de ce dossier, « Groupes de
Barsotti--Tate et cristaux », et rédigé par Messing en 1972.}

\pagerange{18}{19}
\subsection*{La description étale (pages 18 et 19)}

Pour un schéma abélien $A$ \emph{ordinaire} sur une base $S$ de
caractéristique $p$, la décomposition $H^1_{\mathrm{dR}}(A/S) \simeq t_A^\vee
\oplus t_{A^*}$ vaut, et l'on a de plus
\[
t_A \simeq \mathcal{L}_A \otimes_{\mathbb{F}_p} \mathcal{O}_S , \qquad
t_{A^*} \simeq \mathcal{L}_{A^*} \otimes_{\mathbb{F}_p} \mathcal{O}_S ,
\]
où $\mathcal{L}_A$, $\mathcal{L}_{A^*}$ sont des faisceaux étales localement
libres de rang $g$ sur $\mathbb{F}_p$ :
\[
\mathcal{L}_A = (A^*[p]_{\text{ét}})^\vee , \qquad \mathcal{L}_{A^*} =
(A[p]_{\text{ét}})^\vee .
\]
En effet, l'hypothèse fait de $A[p] \to S$ l'extension canonique $0 \to
A[p]^\circ \to A[p] \to A[p]_{\text{ét}} \to 0$, le premier morphisme radiciel
surjectif, le second étale, tous deux finis localement libres ; la partie
connexe $A[p]^\circ$ est de type multiplicatif, duale de Cartier de
$A^*[p]_{\text{ét}}$, et l'algèbre de Lie de $\underline{\mathrm{Hom}}(E,
\mu_p)$ est $E^\vee \otimes \mathcal{O}_S$. Sur un corps algébriquement clos, la
page le redit en cohomologie :
\[
\mathcal{L}_{A^*} = H^1(A, \mathcal{O}_A)^{F} = H^1_{\text{ét}}(A,
\mathbb{Z}/p) = \operatorname{Hom}(A[p](k), \mathbb{Z}/p) ,
\]
la première égalité par la suite d'Artin--Schreier, d'où $t_A^\vee =
A^*[p](k) \otimes_{\mathbb{F}_p} k$ et $t_{A^*} = A[p](k)^\vee
\otimes_{\mathbb{F}_p} k$.\footnote{La page dit « si $S$ est le spectre d'un
corps parfait » ; les égalités entre points ne valent telles quelles que sur
un corps algébriquement clos, ou en lisant les deux membres comme faisceaux
étales.}

Voilà la stratification cherchée : un faisceau étale tensorisé par
$\mathcal{O}_S$ en porte une canonique, puisque les faisceaux étales sont
insensibles aux épaississements nilpotents. C'est, sous une autre forme, ce que
la théorie de Serre--Tate dit des déformations des variétés abéliennes
ordinaires : elles sont gouvernées par leurs parties étales et
multiplicatives, qui ne se déforment pas.\footnote{La théorie de Serre et Tate
(1964) est antérieure et Grothendieck la connaissait ; la page ne la cite
pas, et on n'affirme pas qu'elle s'en inspire.}

\pagerange{19}{21}
\section*{Polarisation et invariant de Hasse (pages 19 à 21)}

\subsection*{La polarisation (page 19)}

Soit $\varphi$ une polarisation de $A$, section de $\underline{\mathrm{NS}}_{A/S}$,
et $\tilde\varphi : A \to A^*$ l'isogénie qu'elle définit. Sa classe de Chern
est une section de $H^2_{\mathrm{dR}} = \bigwedge^2 H^1_{\mathrm{dR}}(A/S)$. Selon
la décomposition ordinaire,
\[
\textstyle\bigwedge^2 H^1_{\mathrm{dR}}(A/S) \simeq \bigwedge^2 t_A^\vee \oplus
(t_A^\vee \otimes t_{A^*}) \oplus \bigwedge^2 t_{A^*} ,
\]
et la classe n'a de composante que dans le terme médian : les deux facteurs sont
isotropes, ce qui revient à dire que les filtrations de Hodge et conjuguée
sont lagrangiennes. La composante médiane est une section de $t_A^\vee \otimes
t_{A^*} = \underline{\mathrm{Hom}}(t_A, t_{A^*})$, et c'est, au signe près,
l'application tangente $\operatorname{Lie}(\tilde\varphi) : t_A \to
t_{A^*}$.\footnote{La page écrit d'abord que la classe « définit un
homomorphisme $t_A^\vee \otimes t_{A^*} \to \mathcal{O}_S$, i.e.\ $t_{A^*} \to
t_A$ », ce qui est le dual de ce qu'il faut : une \emph{section} de $t_A^\vee
\otimes t_{A^*}$ est une flèche $t_A \to t_{A^*}$. L'accolade sous
$\operatorname{Hom}(t_A, t_{A^*})$, et la ligne suivante, ont la bonne
variance, et c'est elle qu'on garde. Tout le passage est en outre traversé de
diagonales, sans doute annulé.} Par la description étale, elle se déduit par
$\otimes_{\mathbb{F}_p} \mathcal{O}_S$ de $\mathcal{L}_A \to \mathcal{L}_{A^*}$,
transposée de $A[p]_{\text{ét}} \to A^*[p]_{\text{ét}}$ induite par
$\tilde\varphi$.

\pagerange{20}{21}
\subsection*{L'invariant de Hasse (pages 20 et 21)}

Soit à nouveau $g$ la dimension relative de $A/S$, maintenant
\emph{quelconque}. L'algèbre de Lie de $A$ est une $p$-algèbre de Lie ;
l'application $p$-linéaire $x \mapsto x^{[p]}$, ou ce qui revient au même le
Verschiebung, donne $V : t_A^{(p)} \to t_A$. En passant au déterminant,
$(\det t_A)^{\otimes p} \to \det t_A$, d'où une section
\[
\sigma_A \in \Gamma(S, (\det t_A)^{\otimes(1-p)}) = \Gamma(S,
\omega_A^{\otimes(p-1)}) .
\]
« Cette section mérite le nom d'\emph{invariant de Hasse} de $A$ » : ses zéros
sont exactement les points $s$ où $A_s$ n'est pas ordinaire.\footnote{La page
écrit $\omega_A^{(p-1)}$ et, juste avant, $L^{\otimes(p-1)}$ : il s'agit de la
puissance tensorielle, non d'un tordu par le Frobenius. Une première approche,
par l'ouvert ordinaire $U$ et la densité, est écrite puis annulée au milieu
de la page 20.} C'est la définition reçue aujourd'hui.

Quel est le rapport entre $\omega_A$ et $\omega_{A^*}$ ? La suite de Hodge
donne $\det t_A^\vee \otimes \det t_{A^*} \simeq \det H^1_{\mathrm{dR}}(A/S)
\simeq H^{2g}_{\mathrm{dR}}(A/S) \simeq \mathcal{O}_S$, la dernière flèche par
la trace ; d'où $\omega_A \otimes \omega_{A^*}^{\vee} \simeq \mathcal{O}_S$,
c'est-à-dire un isomorphisme canonique
\[
\omega_A \simeq \omega_{A^*} .
\]
La page affirme qu'il est compatible avec les « $p$-structures », donc
échange $\sigma_A$ et $\sigma_{A^*}$, et que, pour $A$ ordinaire, il induit un
isomorphisme $\det A[p]_{\text{ét}} \simeq \det A^*[p]_{\text{ét}}$ ---
puisque, par la page 19, $\omega_A = \det(A^*[p]_{\text{ét}}) \otimes
\mathcal{O}_S$ et $\omega_{A^*} = \det(A[p]_{\text{ét}}) \otimes
\mathcal{O}_S$.\footnote{« Je dis que » : l'affirmation n'est pas démontrée,
ni ici ni plus loin, et le mot « $p$-structures » est d'une lecture
incertaine. On la donne pour ce qu'elle est.} C'est l'annonce de ce qui suit.

\pagerange{22}{28}
\section*{L'isomorphisme $\det T_p(A)_{\text{ét}} \simeq \det T_p(A^*)_{\text{ét}}$ (pages 22 à 28)}

\pagerange{22}{23}
\subsection*{Le problème, et le modèle $\ell$-adique}

Soit $A$ ordinaire sur $S$, et $T_p(A)_{\text{ét}} = \varprojlim
A[p^\nu]_{\text{ét}}$, faisceau $p$-adique lisse de rang $g$. Le module de
Tate $p$-adique $T_p(A) = \varprojlim A[p^\nu]$ (comme schéma en groupes) est
extension
\[
0 \to T_p(A)_{\mathrm{rad}} \to T_p(A) \to T_p(A)_{\text{ét}} \to 0 ,
\]
et la dualité de Cartier, compatible aux filtrations, donne
\[
T_p(A)_{\mathrm{rad}} \simeq D(T_p(A^*)_{\text{ét}}) , \qquad
T_p(A^*)_{\mathrm{rad}} \simeq D(T_p(A)_{\text{ét}}) .
\]
Les deux parties étales déterminent donc les deux facteurs de $T_p(A)$, et
$T_p(A)$ lui-même si la base est parfaite, où la suite se scinde
canoniquement. On voudrait un isomorphisme canonique\footnote{La page 23 écrit le but sans l'indice « ét », que la chemise de la
page 14 et toute la suite portent.}
\[
\det T_p(A)_{\text{ét}} \xrightarrow{\ \sim\ } \det T_p(A^*)_{\text{ét}} .
\]

Pour $\ell \neq p$, il existe : $\det T_\ell(A) \simeq \det T_\ell(A^*) \simeq
\mathbb{Z}_\ell(g)$. Et si $\varphi$ est une polarisation, $\beta_\ell(\varphi)
= \det T_\ell(\tilde\varphi)$ donne une flèche $\det T_\ell(A) \to \det
T_\ell(A^*)$, dont l'isomorphisme canonique est $\alpha_\ell =
\beta_\ell(\varphi)/\deg\tilde\varphi$.\footnote{La page écrit « $\deg
\varphi$ » pour $\deg\tilde\varphi$. La division par le degré entier est la
bonne à $\ell$ : $\det T_\ell(u) = \deg u$ pour toute isogénie $u$, rang $2g$
oblige, de sorte que le quotient ne dépend pas de $\varphi$. Une
justification par $T_\ell(A^*) \simeq T_\ell(A)^\vee(1)$ est essayée puis
annulée.} À $p$, le rang n'est plus que $g$, et la normalisation naturelle
devient la racine carrée du degré. On définit donc
\[
\beta_p(\varphi) = \det T_p(\tilde\varphi)_{\text{ét}} , \qquad
\alpha_p(\varphi) = \frac{\beta_p(\varphi)}{\sqrt{\deg\tilde\varphi}} :
\det T_p(A)_{\text{ét}} \otimes \mathbb{Q}_p \to \det T_p(A^*)_{\text{ét}}
\otimes \mathbb{Q}_p ,
\]
le degré d'une polarisation étant un carré parfait.\footnote{Le second
$\otimes\mathbb{Q}$ porte l'indice $\ell$ sur la page : un lapsus pour $p$.
Le degré de $\tilde\varphi$ est le carré de la caractéristique d'Euler du
fibré en droites (Riemann--Roch) ; $\sqrt{\deg\tilde\varphi}$ est sa racine
positive.} Il reste à prouver : a) $\alpha_p(\varphi)$ ne dépend pas de
$\varphi$ ; b) c'est un isomorphisme des réseaux, non seulement des
$\mathbb{Q}_p$-droites.

On suppose ici que $A$ est polarisable, au moins localement sur $S$, et l'on
raisonne sur une base connexe, où les déterminants d'endomorphismes de
faisceaux $p$-adiques de rang $g$ sont des scalaires de $\mathbb{Z}_p$.\footnote{Hypothèses
implicites sur la page. La première est ce qui permet de recoller les
$\alpha_p(\varphi)$ locaux, une fois a) démontré.}

Un lemme vient ensuite, qui affirmait que pour tout endomorphisme $u$ de $A$,
$\deg u$ est un carré et $\det T_p(u)_{\text{ét}} = \sqrt{\deg u}$. Il est
barré de quatre traits et marqué « faux » dans la marge.\footnote{Il l'est en
effet. Sur un corps fini $\mathbb{F}_q$, pour $u$ l'endomorphisme de Frobenius
d'une variété ordinaire, $\det T_p(u)_{\text{ét}}$ est le produit des $g$
racines unités du polynôme caractéristique, une unité $p$-adique, alors que
$\deg u = q^g$. Le contre-exemple est le nôtre.}

\pagerange{24}{26}
\subsection*{a) L'indépendance de la polarisation (pages 24 à 26)}

Soit $\psi$ une autre polarisation : $\tilde\psi = \tilde\varphi \circ u$ avec
$u \in \operatorname{End}(A) \otimes \mathbb{Q}$. La symétrie de $\tilde\psi$
et de $\tilde\varphi$ donne $u^* \tilde\varphi = \tilde\varphi u$, soit
\[
u^* = \tilde\varphi\, u\, \tilde\varphi^{-1} ,
\]
c'est-à-dire que $u$ est symétrique pour l'involution de Rosati $u \mapsto
u^\dagger = \tilde\varphi^{-1} u^* \tilde\varphi$ ; et $u$ est même
\emph{positif}, comme quotient de deux polarisations.\footnote{La page écrit
« $u' = \tilde\varphi u \tilde\varphi^{-1}$ » puis appelle $u \mapsto u'$
l'involution de $\operatorname{End}(A) \otimes \mathbb{Q}$. Mais
$\tilde\varphi u \tilde\varphi^{-1}$ est un endomorphisme de $A^*$, non de
$A$ : c'est $u^*$ quand $u$ est symétrique, et c'est sous ce nom que la
page l'utilise ensuite. L'involution elle-même est celle de Rosati,
$\tilde\varphi^{-1} u^* \tilde\varphi$. Le signe de « $\tilde\varphi' =
-\tilde\varphi$ » tient à sa convention pour la bidualité et disparaît dans
la conclusion.} Posons $\Delta(u) = \det T_p(u)_{\text{ét}}$ (étendu à
$\operatorname{End}(A) \otimes \mathbb{Q}$, à valeurs dans $\mathbb{Q}_p$) et
$\delta(u) = \sqrt{\deg u}$, racine positive --- $\deg u = \deg\tilde\psi /
\deg\tilde\varphi$ est le carré d'un rationnel. Comme $\beta_p(\psi) =
\beta_p(\varphi)\Delta(u)$ et $\sqrt{\deg\tilde\psi} =
\sqrt{\deg\tilde\varphi}\,\delta(u)$, l'énoncé a) équivaut à
\[
\Delta(u) = \delta(u) \qquad \text{pour tout } u \text{ symétrique positif.}
\]

La preuve a deux temps. D'abord, le carré :
\[
\Delta(u)\,\Delta(u^*) = \deg u = \delta(u)^2 ,
\]
et comme $u^* = \tilde\varphi u \tilde\varphi^{-1}$, $\Delta(u^*) = \Delta(u)$,
d'où $\Delta(u)^2 = \delta(u)^2$ et $\Delta(u) = \pm\delta(u)$.\footnote{La
première égalité est introduite par « Or par \ldots\ on a », le nom de la
justification étant illisible ; la page 25, annulée, l'attribuait au
« théorème $p$-adique ou th.\ caractéristique de Weil ». Elle est vraie, et
voici pourquoi, de notre fait : pour $A$ ordinaire sur un corps parfait, le
module de Dieudonné de $A[p^\infty]$ est somme de sa partie étale et de sa
partie multiplicative ; $u$ opère sur la première avec le déterminant
$\Delta(u)$, sur la seconde --- duale de Cartier de $A^*[p^\infty]_{\text{ét}}$
--- avec le déterminant $\Delta(u^*)$ ; et le déterminant de $u$ sur tout le
module de Dieudonné est $\deg u$, terme constant du polynôme caractéristique
de Weil. Le dossier n'établit nulle part cette identité ; il l'utilise.}

Ensuite, le signe. Sur le sous-espace des éléments symétriques de
$\operatorname{End}(A) \otimes \mathbb{Q}$, $\Delta$ est une fonction
polynomiale en les coordonnées, à coefficients dans $\mathbb{Z}_p$, qui prend
des valeurs rationnelles aux points rationnels (puisqu'elle y vaut $\pm\delta$) :
ses coefficients sont donc rationnels, et elle s'étend aux éléments
symétriques de $\operatorname{End}(A) \otimes \mathbb{R}$, à valeurs réelles et
multiplicative. Or dans cette algèbre à involution positive, tout élément
symétrique positif est le carré $v^2$ d'un élément symétrique $v$, et
$\Delta(v^2) = \Delta(v)^2 \geq 0$. Donc $\Delta \geq 0$ sur le cône positif,
et $\Delta(u) = +\delta(u)$. Ce qu'il fallait démontrer.\footnote{La page
affirme la rationalité des coefficients de $\Delta$ sans restreindre aux
éléments symétriques ; sur tout $\operatorname{End}(A)$, $\Delta$ prend des
valeurs $p$-adiques non rationnelles (cf.\ le Frobenius ci-dessus) et
l'argument ne passe pas. Restreint aux symétriques, il est exact, et c'est
tout ce que demande la positivité. La page 25, annulée, essayait une autre
voie, par un caractère de signe $\varepsilon : F^\times \to \{\pm 1\}$ sur $F =
\operatorname{End}(A) \otimes \mathbb{Q}$.}

\pagerange{26}{27}
\subsection*{b) Un isomorphisme des réseaux (pages 26 et 27)}

$\alpha_p(\varphi)$ envoie $\det T_p(A)_{\text{ét}}$ \emph{sur} $\det
T_p(A^*)_{\text{ét}}$ si et seulement si
\[
v_p\bigl(\operatorname{card}\operatorname{Coker} T_p(\tilde\varphi)_{\text{ét}}\bigr)
= v_p\bigl(\sqrt{\deg\tilde\varphi}\bigr) .
\]
Or $\operatorname{Coker} T_p(\tilde\varphi)_{\text{ét}} \simeq (\operatorname{Ker}
\tilde\varphi)(p)_{\text{ét}}$, partie étale de la composante $p$-primaire du
noyau. Le noyau d'une polarisation est autodual ; sa partie étale est donc
duale de sa partie multiplicative, et, $A$ étant ordinaire, il n'a pas de
partie locale-locale. D'où
\[
\bigl(\operatorname{card}(\operatorname{Ker}\tilde\varphi)(p)_{\text{ét}}\bigr)^2
= \operatorname{card}(\operatorname{Ker}\tilde\varphi)(p) =
(\deg\tilde\varphi)_{(p)} ,
\]
la $p$-partie du degré, ce qui est la relation voulue. On obtient un
isomorphisme canonique, indépendant de toute polarisation,
\[
\alpha_A : \det T_p(A)_{\text{ét}} \xrightarrow{\ \sim\ } \det
T_p(A^*)_{\text{ét}} ,
\]
compatible au changement de base et naturel pour les isomorphismes de schémas
abéliens.

\pagerange{27}{28}
\subsection*{Fonctorialité, et ce qui reste à vérifier (pages 27 et 28)}

La naturalité pour les isomorphismes dit : si $u$ est un automorphisme de $A$,
et $\check u = (u^*)^{-1}$ l'automorphisme qu'il induit sur $A^*$, alors
$\alpha_A$ entrelace $\det T_p(u)_{\text{ét}}$ et $\det T_p(\check
u)_{\text{ét}}$ ; en identifiant les deux déterminants à $\mathbb{Z}_p$,
\[
\det T_p(u)_{\text{ét}} \cdot \det T_p(u^*)_{\text{ét}} = 1 .
\]
Pour un endomorphisme quelconque, en revanche, $\det T_p(u)_{\text{ét}} \neq
\det T_p(u^*)_{\text{ét}}$ en général, et l'on ne peut demander à $\alpha$
d'être naturel pour les isogénies : $A \mapsto \det T_p(A)_{\text{ét}}$ est
covariant, $A \mapsto \det T_p(A^*)_{\text{ét}}$ contravariant --- la page le
dit elle-même, « les deux foncteurs n'étant pas de même variance en $A$ ».
Ce qui tient lieu de fonctorialité est la formule encadrée
\[
\det T_p(u)_{\text{ét}} \cdot \det T_p(u^*)_{\text{ét}} = \deg u .
\]
Elle se démontre dans le cadre ci-dessus, et sous une forme plus générale : pour
une isogénie $f : A \to B$ de schémas abéliens ordinaires polarisables,
\[
\det T_p(f^*)_{\text{ét}} \circ \alpha_B \circ \det T_p(f)_{\text{ét}} = (\deg f)\,
\alpha_A ,
\]
en appliquant a) à la polarisation $f^*\tilde\varphi_B f$ de $A$.\footnote{Cette
forme, et la dérivation à partir de a), sont les nôtres ; la page encadre le
cas $B = A$ et l'étend, dans une marge très incertaine, aux isogénies entre
schémas abéliens de même dimension. Pour $B = A$ on retrouve l'identité de la
page 24, qui en est en réalité l'\emph{entrée} : la dérivation ne la démontre
donc pas, elle en montre la cohérence avec a).}

Deux vérifications sont laissées ouvertes :
\begin{itemize}
\item[c)] $\alpha_{A^*} = \alpha_A^{-1}$, « signe (?) ». Si $\tilde\varphi$ est
une polarisation de $A$ et $n$ tue son noyau, $n\tilde\varphi^{-1}$ est une
polarisation de $A^*$, et le calcul direct donne $\alpha_{A^*} =
\alpha_A^{-1}$ --- au signe près qu'introduit la convention d'identification
$A^{**} = A$, qui est précisément ce que la page met en doute.\footnote{Ce
calcul est le nôtre. La phrase suivante de la page, « Noter aussi la
compatibilité par \ldots\ de $\mathcal{V}A$ », est trop lacunaire pour être
rendue.}
\item[d)] L'isomorphisme $\det A[p]_{\text{ét}} \simeq \det
A^*[p]_{\text{ét}}$ induit par $\alpha_A$ est celui que la dualité de De Rham
donnait à la page 21 --- « à vérifier ». Ce ne l'est pas dans le dossier, et
la boucle entre l'invariant de Hasse et l'isomorphisme $p$-adique reste
ouverte là.\footnote{La page écrit « l'isom.\ $({}_pA)_{\text{ét}} \to
({}_pA^*)_{\text{ét}}$ », sans « det » ; à la page 21 il s'agit des
déterminants, et c'est ce que l'on rétablit.}
\end{itemize}

\end{document}
